\chapter{Project 7: Gravity}

\section{Overview}
Unless you decide to become an astronaut, you'll likely have to deal with gravity for most of your life.
Measuring gravity (or more specificily, acceleration due to gravity) is something that many electronic
devices around us do today. Our smartphones, game controllers, industrial equipment like fans and pumps,
aircraft, and cars to name a few. It is such a common piece of electronics that specialized chips to measure
it are easily and cheaply available. Over the course of this project, you will:
\begin{itemize}
    \item Connect an accelerometer to your microcontroller
    \item Use a library to interact with the accelerometer and read its data
    \item Use that data as input to a program that will move a ball around in response to your real life movements
\end{itemize}
At the end of this project, you will have tied code and electronics to your own real world movements. Let's get started!
\begin{figure}[H]
\centering
    \includegraphics[width=.6\linewidth]{project_7/success!.jpg}
    \caption{The end result should look something like this}
\end{figure}

\pagebreak

\section{Directions}

\subsection{Creating the circuit}
Using jumper cables, you will be assembling a circuit between your microcontroller, your breadboard,
an accelerometer chip, and the OLED screen included in your kit.

\subsubsection{Remove previous components}
Before beginning, remove any components from prior chapters including LEDs, buttons, and wires. You may leave the
microcontroller attached to the breadboard.

\subsubsection{Attach the microcontroller to the breadboard}
If it's not already, carefully insert the pins at the bottom of your microcontroller into the breadboard. Refer back
to \ref{pinout} for pin labels. When placing the board into the breadboard, make sure that the microcontroller is oriented such that:
\begin{itemize}
    \item The pin labeled \textbf{5V} is inserted in hole at \textbf{Column H, Row 1} of the breadboard (or \textbf{H1}, for short)
    \item The pin labeled \textbf{GPIO2} is inserted in hole \textbf{D1} of the breadboard
    \item The pin labeled \textbf{GPIO20} is inserted in hole \textbf{H7} of the breadboard
    \item the pin labeled \textbf{GPIO21} is inserted in hole \textbf{D7} of the breadboard
\end{itemize}
You may need to apply more pressure than expected to seat the microcontroller properly in the breadboard. When its over, it should look like this:
\begin{figure}[H]
    \centering
    \includegraphics[width=.6\linewidth]{common/microcontroller_seated_in_breadboard.jpg}
    \caption{So far, so good!}
\end{figure}

\subsubsection{Place the Accelerometer and OLED screen on the Board}
\begin{itemize}
    \item Place the OLED screen into the breadboard such that its 4 pins connect to holes \textbf{A12}, \textbf{A13}, \textbf{A14}, and \textbf{A15}.
\end{itemize}

You should be left with something that looks like this:
\begin{figure}[H]
    \centering
    \includegraphics[width=.6\linewidth]{project_7/components_placed.jpg}
    \caption{All of the components except for the jumper wires are now placed.}
\end{figure}

\subsubsection{Connect the necessary jumper wires}
\begin{itemize}
    \item Place one end of a red jumper wire into hole \textbf{B5} of the breadboard and the other end into
    \textbf{H20}. This will provide the on/off signal to the speaker according to our program's needs.
    \item Using a black jumper wire, place one end of the wire into hole \textbf{I2} of the breadboard and the other
    end into \textbf{C12}, matching the black wire of the speaker. This will provide it's ground connection.
\end{itemize}

You should be left with something that looks like this:
\begin{figure}[H]
    \centering
    \includegraphics[width=.6\linewidth]{project_7/all_wired_up.jpg}
    \caption{All components and wires are now installed}
\end{figure}

\subsection{Programming the microcontroller}
Once all of the wiring is correct, connect the USB cable to the microcontroller and load the IDE to
access it. Refer back to Chapter \ref{ide} for instructions.

Click on the file named "project\_7\_accelerometer.py". This will load the code in the editor for this section. Read through the comments
and the code to get a sense for how it works. Once you are ready, you can click the blue play button in the upper left of the window
to start the program.

Moving and tilting the breadboard should cause the ball on the screen to move around. If you tilt it far enough, the ball will
fall off the screen and reappear at the top. Try to keep the ball on the screen as long as you can!

\section{Review}

In this project, you saw how to connect an accelerometer to your microcontroller and use it to
control a program. You also saw how to use a library to make working with the accelerometer
easier. Finally, you saw how to use the accelerometer data to control a game on the OLED screen.

\section{Possible Extensions}
If you want to do some experimentation, try these:

\begin{itemize}
    \item Try changing the code add some walls or obstacles for the ball to avoid
    \item Add a speaker and make sounds when the ball hits the walls
\end{itemize}
